 \section{行き詰まったとき}\label{節：解法の変更}

  変形で簡単にならなければ,残った項を次の手掛かりにする.例えば,
  \begin{enumerate}
    \item \bookterm{logarithm}{対数}を取っても,\bookterm{sequence}{数列}の項を1次の和として表せない.
    \item 定数を引いて付加項を消そうとしたが,条件を満たす定数が求まらない.
    \item 置換後に新しい種類の項が増え続ける.
  \end{enumerate}
  定数で消せない$n$の項なら一次式を引く.対数で和を処理できなければ,先に因数分解や定数の調整を試す.
  \BookFigFallback

  $a_2$から$a_5$程度の実験も予想に役立つが,\bookterm{initial}{初期条件}と更新式で証明する必要がある.
  一つの方法がうまくいかなくても,その漸化式が解けないと決まったわけではない.別の公式や初項の値にも注目してみよう.

  \bookterm{general}{一般項}が難しくても,正値・単調性・\bookterm{limit}{極限}・保存される量は調べられるかもしれない.
  \bookproblemref{practice-entrance}{6}{発展問6}では逆数の更新から減り方を求める.
  求めたい性質に合わせて,追う量と目標を選び直す.
